\newcommand{\ProteinDataCaption}{Original graph-digitized Tang et al. (2018) 100 mL isothermal lysozyme observations at 20 degrees C. (a) Dissolved concentration with assumed 5\% standard deviations; the diamond fixes the initial condition, filled circles train and open squares test. (b) Training conditional number densities at 74, 180 and 420 min. (c) Held-out 660-min density. (d) Conditional number $d_{10},d_{50},d_{90}$; connecting dotted lines are guides only. All sizes and quantiles condition on 2.5--25 $\mu$m. No model prediction or invented PSD uncertainty is shown.}
\newcommand{\ProteinFitDetails}{(a) Concentration; (b) training conditional number densities at 74, 180 and 420 min; (c) held-out 660-min density; (d) conditional number quantiles. All densities and quantiles condition on 2.5--25 $\mu$m. Solid curves/stairs show the inference-coordinate MAP; observations are dashed density reconstructions and markers. Shading is the pointwise 95\% latent interval from screened NUTS draws, conditional on fixed observation errors and physical constants; it excludes additional measurement noise. The vertical dotted line marks the end of training at 420 min. The 660-min PSD and later concentrations are predicted without refitting. Quantile segments are omitted when the conditioning window is empty; quantiles derived from PSDs are not independent likelihood data. Numerical densities use actual bin edges without smoothing.}
\newcommand{\ProteinCornerDetails}{Diagonals show normalized posterior KDEs and dashed analytical priors generated from the active specifications; positive wide-scale parameters use labelled natural-log plotting coordinates, including the correct prior Jacobian. Broad priors are not renormalized to the displayed posterior window; full support is in the prior table. Where shown, joint KDE contours enclose approximately 68\% and 95\% probability. The thin, curved nucleation-prefactor/interfacial-energy relation is shown directly with retained draws. Orange markers denote the transformed inference-coordinate MAP. Sampling convergence does not establish numerical convergence, identifiability or physical model adequacy.}
\newcommand{\ProteinCNTFitCaption}{CNT nucleation and size-independent mean growth. \ProteinFitDetails MUSCL growth and native SSPRK33 use 400 bins on 0--50 $\mu$m with combined CFL 0.35 and maximum step 2 min; 200 retained parameter draws supply the bands. The posterior retains 4 chains of 500 draws each.}
\newcommand{\ProteinCNTCornerCaption}{CNT nucleation and size-independent mean growth: NUTS posterior. \ProteinCornerDetails}
\newcommand{\ProteinGRDFitCaption}{CNT nucleation and mean growth with random-fluctuation dispersion $D_G=\theta_{\mathrm{GRD}}G^2$. \ProteinFitDetails MUSCL growth and native SSPRK33 use 100 bins on 0--50 $\mu$m with combined CFL 0.35 and maximum step 2 min; 200 retained parameter draws supply the bands. The posterior retains 4 chains of 2000 draws each.}
\newcommand{\ProteinGRDCornerCaption}{CNT nucleation and mean growth with random-fluctuation dispersion $D_G=\theta_{\mathrm{GRD}}G^2$: NUTS posterior. \ProteinCornerDetails}
\newcommand{\ProteinPriorTableCaption}{Implemented priors, physical units and applicable models. Inference transformations are separate from prior-density coordinates; the dispersion Exponential prior has mean/scale 10 min.}
\newcommand{\ProteinParameterTableCaption}{Inference-coordinate MAP estimates transformed to physical units, posterior medians and equal-tailed 95\% credible intervals. Model 1 fixes dispersion to zero. Intervals condition on the model, fixed observation errors and physical constants.}
